\subsection{Архимед. Сложный}
%%%%%%%%%%%%%%%%%%%%%%%%%%%%%%%%%%%%%%%%%%%%%%%%%%%%
%%%%%%%%%%%%%%%%%%%%%%%%%%%%%%%%%%%%%%%%%%%%%%%%%%%%
 \z[{\begin{wrapfigure}[5]{r}{0.15\textwidth}%
		\vspace{-15pt}%
		\begin{flushright}%
 			\includegraphics[width=0.15\textwidth]{9.3.1}%
		\end{flushright}%
	\end{wrapfigure}%
	}]%
	{Два кубика, связанные натянутой нитью, находятся в воде. Верхний кубик со стороной $a = 10\,\text{см}$ плавает, погрузившись в воду на три четверти своего объёма. Нижний кубик касается дна (вода под него подтекает). Сторона нижнего кубика равна $a/2$, а его плотность в 2 раза больше, чем у верхнего. Определите, при каких значениях плотности $\rho$ материала верхнего кубика возможно такое состояние системы.} 
%
{$700\,\text{кг/м}^3 < \rho < 750\,\text{кг/м}^3$}
%%%%%%%%%%%%%%%%%%%%%%%%%%%%%%%%%%%%%%%%%%%%%%%%%%%%
\z[{\begin{wrapfigure}[5]{r}{0.15\textwidth}%
		\vspace{-15pt}%
		\begin{flushright}%
 			\includegraphics[width=0.15\textwidth]{9.3.2}%
		\end{flushright}%
	\end{wrapfigure}%
	}]%
	{Деревянный шарик привязан ко дну цилиндрического сосуда с водой. Над поверхностью воды находится часть шарика, а нить натянута. Если нить перерезать, то шарик всплывёт, и уровень воды в сосуде изменится на $\Delta h = 4\,\text{см}$. С какой силой $F$ была натянута нить? Площадь дна сосуда $S = 100\,\text{см}^2$.}
%
{$T = 4\,\text{Н}$}
%%%%%%%%%%%%%%%%%%%%%%%%%%%%%%%%%%%%%%%%%%%%%%%%%%%% 
 \z[{\begin{wrapfigure}[9]{r}{0.2\textwidth}%
		\vspace{-15pt}%
		\begin{flushright}%
 			\includegraphics[width=0.2\textwidth]{9.3.3}%
		\end{flushright}%
	\end{wrapfigure}%
	}]%
	{Водолазный колокол в форме цилиндра без дна, частично заполненный воздухом, находится под водой. Чтобы колокол не всплывал, его прикрепили тросом к дну водоёма. На верёвке к колоколу привязан груз, находящийся в воде. Площадь сечения колокола $S = 4\,\text{м}^2$, объём воздуха в нём $V = 8\,\text{м}^3$ при давлении $p = 1{,}5\cdot10^5\,\text{Па}$. Когда груз в колоколе поднимают над уровнем воды, давление возрастает на $\Delta p = 250\,\text{Па}$, при этом трос остаётся натянутым. Найдите изменение натяжения $\Delta T$ троса и верёвки. Воздух в колоколе подчиняется закону Бойля-Мариотта: $pV = \text{const}$, где $p$ -- давление, $V$ -- объём воздуха.}
%
{$\Delta T_{\text{троса}} = \dfrac{\rho g V \Delta p}{p + \Delta p} = 133\,\text{Н}$; ~ $\Delta T_{\text{веревки}} = \Delta T_{\text{троса}} + S\Delta p = 1133\,\text{Н}$.}
%%%%%%%%%%%%%%%%%%%%%%%%%%%%%%%%%%%%%%%%%%%%%%%%%%%% 
 \z[{\begin{wrapfigure}[4]{r}{0.2\textwidth}%
		\vspace{-15pt}%
		\begin{flushright}%
 			\includegraphics[width=0.2\textwidth]{9.3.4}%
		\end{flushright}%
	\end{wrapfigure}%
	}]%
	{Подъёмный кран медленно поднимает с помощью троса плавающее в воде бревно (рис.). Трос прикреплён к одному концу бревна, которое можно считать тонким цилиндром с постоянной плотностью. Масса бревна $m$, длина — $L$. Отношение плотностей воды и древесины $\gamma = \frac{4}{3}$. Ускорение свободного падения $g$.
	\lpu Какую минимальную работу $A_{\text{min}}$ нужно совершить крану, чтобы полностью вытащить бревно из воды? \lpuend
	\lpu Постройте график зависимости силы натяжения $T$ троса от высоты над водой $h$ приподнимаемого конца бревна. Укажите характерные точки графика. \lpuend
	\lpu Какую работу $A_h$ совершит кран при переводе бревна из одного наклонного положения в другое наклонное положение, в котором верхний конец бревна поднялся на высоту $\Delta h = \frac{L}{5}$? \lpuend
	}
%
{\puan $A_{\text{мин}} = \dfrac{1}{2}mgL$;~\puan $T = \dfrac{mg}{3}$, если $0 \leq h \leq \dfrac{L}{2}$;~  $T = \dfrac{mg}{3}\left(\dfrac{4h}{L} - 1\right)$, если $\dfrac{L}{2} \leq h \leq L$;~ $T = mg$, если $L \leq h$;~\puan $A_h = \dfrac{1}{15}mgL$}
%%%%%%%%%%%%%%%%%%%%%%%%%%%%%%%%%%%%%%%%%%%%%%%%%%%%	
 \z[{\begin{wrapfigure}[3]{r}{0.2\textwidth}%
		\vspace{-15pt}%
		\begin{flushright}%
 			\includegraphics[width=0.2\textwidth]{9.3.5}%
		\end{flushright}%
	\end{wrapfigure}%
	}]%
	{Определите минимальную силу натяжения $F$ двух канатов, связывающих широкий плот, состоящий из двух слоев бревен. Масса каждого бревна $m$. Верхний слой бревен погружен в воду наполовину.}
%
{$T = \dfrac{\sqrt{3}mg}{72}$}